% !TEX program = xelatex
% =============================================================================
%  《模拟滤波器设计：透过现象看本质（巴特沃斯与切比雪夫）》
%  深度学习笔记与刷题指南
%  作者：资深 DSP 架构师
% =============================================================================

\documentclass[12pt,a4paper]{ctexart}

% ------------------ 页面布局 ------------------
\usepackage[top=2.5cm, bottom=2.5cm, left=2.8cm, right=2.8cm]{geometry}
\usepackage{setspace}
\onehalfspacing

% ------------------ 数学与排版 ------------------
\usepackage{amsmath,amssymb,amsfonts}
\usepackage{mathtools}
\usepackage{bm}
\usepackage{upgreek}
\usepackage{mathrsfs}

% ------------------ 图形与颜色 ------------------
\usepackage{tikz}
\usepackage{tikz-3dplot}
\usetikzlibrary{calc, patterns, decorations.pathmorphing, decorations.markings,
                arrows.meta, shapes.geometric, positioning, backgrounds,
                intersections, fit}

% ------------------ 高亮框 ------------------
\usepackage[framemethod=TikZ]{mdframed}
\usepackage{tcolorbox}
\tcbuselibrary{skins, breakable, most}

% ------------------ 表格 ------------------
\usepackage{array}
\usepackage{booktabs}
\usepackage{colortbl}
\usepackage{longtable}

% ------------------ 列表 ------------------
\usepackage{enumitem}
\setlist{nosep, leftmargin=*}

% ------------------ 超链接与引用 ------------------
\usepackage[hidelinks]{hyperref}
\usepackage{cleveref}

% ------------------ 其他 ------------------
\usepackage{caption}
\usepackage{subcaption}
\usepackage{fancyhdr}
\usepackage{lastpage}
\usepackage{ifthen}

% ==================== 颜色定义 ====================
\definecolor{butterworth}{RGB}{0, 102, 204}
\definecolor{chebyshev}{RGB}{204, 51, 0}
\definecolor{alertred}{RGB}{180, 0, 0}
\definecolor{deepgreen}{RGB}{0, 110, 60}
\definecolor{boxbg}{RGB}{248, 248, 252}
\definecolor{notebg}{RGB}{255, 252, 240}

% ==================== 自定义框环境 ====================

% ---- 4步设计法主框架 ----
\newtcolorbox[auto counter, number within=section]{
    designflow
}[1][]{
    enhanced,
    colback=boxbg,
    colframe=black!60,
    coltitle=white,
    fonttitle=\bfseries\Large,
    title={\hspace{0.5em} 终极工程模板：模拟滤波器 4 步傻瓜设计法},
    attach boxed title to top left={xshift=5mm, yshift=-3mm},
    boxed title style={
        size=small, colback=black!70, sharp corners, boxrule=0pt
    },
    sharp corners,
    before skip=15pt,
    after skip=15pt,
    left=8pt, right=8pt, top=10pt, bottom=8pt,
    #1
}

% ---- 步骤块 ----
\newtcolorbox{stepbox}[1]{
    enhanced,
    colback=white,
    colframe=#1,
    coltitle=white,
    fonttitle=\bfseries\large,
    title={},
    attach boxed title to top left={xshift=3mm, yshift=-2.5mm},
    boxed title style={
        size=small, colback=#1, sharp corners, boxrule=0pt
    },
    sharp corners,
    left=4pt, right=4pt, top=6pt, bottom=4pt,
    before skip=6pt, after skip=6pt,
}

% ---- 避坑提示 ----
\newtcolorbox{pitfall}{
    enhanced,
    colback=notebg,
    colframe=alertred!70,
    fontupper=\small\kaishu,
    title={\textbf{!! 避坑提示}},
    attach boxed title to top left={xshift=3mm, yshift=-2.5mm},
    boxed title style={
        size=small, colback=alertred!80, coltitle=white,
        fontupper=\bfseries\small
    },
    sharp corners,
    left=6pt, right=6pt, top=8pt, bottom=5pt,
}

% ---- 关键洞察框 ----
\newtcolorbox{insight}{
    enhanced,
    colback=deepgreen!5,
    colframe=deepgreen!70,
    fontupper=\kaishu,
    title={\textbf{>> 本质洞察}},
    attach boxed title to top left={xshift=3mm, yshift=-2.5mm},
    boxed title style={
        size=small, colback=deepgreen!80, coltitle=white,
        fontupper=\bfseries\small
    },
    sharp corners,
    breakable,
    left=6pt, right=6pt, top=8pt, bottom=5pt,
}

% ---- 刷题答案框 ----
\newtcolorbox{answerbox}{
    enhanced,
    colback=black!3,
    colframe=black!30,
    fontupper=\small,
    title={\textbf{参考答案}},
    attach boxed title to top left={xshift=3mm, yshift=-2.5mm},
    boxed title style={
        size=small, colback=black!70, coltitle=white,
        fontupper=\bfseries\small
    },
    sharp corners,
    breakable,
    left=8pt, right=8pt, top=8pt, bottom=5pt,
}

% ==================== 页眉页脚 ====================
\pagestyle{fancy}
\setlength{\headheight}{14pt}
\fancyhf{}
\fancyhead[L]{\small\color{gray}模拟滤波器设计：透过现象看本质}
\fancyhead[R]{\small\color{gray}巴特沃斯 \& 切比雪夫}
\fancyfoot[C]{\small\color{gray}---\ \thepage\ ---}
\renewcommand{\headrulewidth}{0.4pt}
\renewcommand{\footrulewidth}{0pt}

% ==================== 常用数学快捷命令 ====================
\newcommand{\Omegac}{\Omega_c}
\newcommand{\Omegap}{\Omega_p}
\newcommand{\Omegas}{\Omega_s}
\newcommand{\db}{\ \mathrm{dB}}
\newcommand{\Hz}{\ \mathrm{Hz}}
\newcommand{\rads}{\ \mathrm{rad/s}}
\newcommand{\splane}{s\text{-plane}}
\newcommand{\Han}{H_{\!a}(s)}
\newcommand{\HanNorm}{H_{\!an}(s)}
\newcommand{\absSq}[1]{\bigl|#1\bigr|^2}
\newcommand{\ch}{\operatorname{ch}}
\newcommand{\arch}{\operatorname{arch}}
\newcommand{\ceil}[1]{\bigl\lceil #1 \bigr\rceil}
\newcommand{\floor}[1]{\bigl\lfloor #1 \bigr\rfloor}
\DeclareMathOperator{\arctanh}{arctanh}
\DeclareMathOperator{\sh}{sh}
\DeclareMathOperator{\arsh}{arsh}

% =============================================================================
\begin{document}

% =========================== 标题页 ===========================
\begin{center}
    {\Huge\bfseries 模拟滤波器设计\\[6pt]}
    {\huge\bfseries 透过现象看本质\\[8pt]}
    {\Large —— 巴特沃斯与切比雪夫 ——\\[15pt]}
    \rule{0.7\textwidth}{1.5pt}\\[10pt]
    {\large\kaishu
    深度学习笔记 \& 刷题指南\\[5pt]
    「公式不是用来背的，是用来理解的」\\[20pt]
    \normalsize\songti 版本 1.0 \quad \textbullet\quad \today
    }
\end{center}

\thispagestyle{empty}
\newpage

% =========================== 目录 ===========================
\setcounter{page}{1}
\tableofcontents
\newpage

% =============================================================================
%  PART I: 核心理论
% =============================================================================
\section{核心理论：两种滤波器的物理本质}

\subsection{从「甲方需求」说起——滤波器到底在解决什么问题？}

想象你坐在工位上，桌面上摊着一份来自甲方（射频前端组）的规格书：

\begin{center}
\begin{tcolorbox}[width=0.85\textwidth, enhanced, colback=white, colframe=black!50,
    sharp corners, halign=center]
    \textbf{低通滤波器技术规格}\\[6pt]
    \begin{tabular}{c|c}
        \toprule
        通带截止频率 $f_p$ & 通带最大衰减 $\alpha_p$ \\
        \midrule
        阻带截止频率 $f_s$ & 阻带最小衰减 $\alpha_s$ \\
        \bottomrule
    \end{tabular}\\[6pt]
    \small\kaishu “我不管你怎么设计，通带里衰减不许超过 $\alpha_p$，阻带里衰减必须大于 $\alpha_s$。超了就别来见我。”
\end{tcolorbox}
\end{center}

这幅图景就是所有滤波器设计的逻辑起点。甲方划出的两条线——通带衰减上限 $\alpha_p$ 和阻带衰减下限 $\alpha_s$——就是\textbf{违约禁区}。你的任务，就是找到刚好满足这些约束的传递函数 $H(s)$，并且用最少的运放来实现它。

\begin{insight}
\textbf{工程经济学的第一性原理：}
\begin{itemize}
    \item \textbf{阶数 $N$ = 成本。} 每一个 $N$ 的实现都对应一个运放、一组 $RC$、一片 PCB 面积和 $N$ 倍功耗。
    \item \textbf{指标余量 = 浪费。} 如果 $\alpha_s$ 只要 $40\db$ 你却做了 $60\db$，你多花的运放并没有人买单。
    \item \textbf{滤波器设计 = 刚好不违约的前提下，阶数最小化。}
\end{itemize}
这就是为什么我们要用巴特沃斯和切比雪夫两种策略——它们代表了「用不同的方式花钱（阶数），买到不同品质的过渡带」。
\end{insight}

% ---------------------------------------------------------------------------
\subsection{幅度平方函数：看清公式背后的逻辑}

模拟滤波器的设计不从传递函数 $H(s)$ 直接下手，因为 $H(s)$ 同时涉及幅度和相位。工程上我们先闭上一只眼——只关心幅度，用\textbf{幅度平方函数} $|H(j\Omega)|^2$ 来干活。

为什么要搞「平方」？因为 $|H(j\Omega)|^2 = H(j\Omega)H(-j\Omega)$，这个函数只含 $\Omega$ 的偶次幂，数学上极其干净。而当我们算出 $|H(j\Omega)|^2$ 之后，再把它「开方」——也就是做谱分解——就能恢复出稳定的 $H(s)$。

% --- 巴特沃斯 ---
\subsubsection{巴特沃斯幅度平方函数：最大平坦的代价}

$$
\boxed{\absSq{H(j\Omega)} = \frac{1}{1 + \left(\dfrac{\Omega}{\Omegac}\right)^{2N}}}
\qquad (N = 1,2,3,\dots)
$$

\textbf{为什么长这样？逐项拆解：}

\begin{enumerate}[leftmargin=2em]
    \item \textbf{分母的 $1 + (\cdots)$ 结构：}
    当 $\Omega = 0$（直流），$|H(j0)|^2 = 1/(1+0) = 1 = 0\db$。这个「$1$」就是用来\textbf{锁死直流增益为 $0\db$} 的。没有它，信号一进来就先被砍一刀，谁也受不了。

    \item \textbf{$2N$ 次方（为什么是偶数次？）：}
    物理上，我们需要 $|H(j\Omega)|^2$ 是 $\Omega$ 的偶函数（幅度响应关于 $\Omega=0$ 对称）。数学上取 $2N$ 保证分母永远 $\ge 1$，函数值不会爆炸。$N$ 每增加 $1$，过渡带额外贡献 $-20\db/\text{dec}$ 的滚降。

    \item \textbf{$\Omegac$ 的角色：}
    当 $\Omega = \Omegac$，$|H|^2 = 1/2$，幅度衰减 $3\db$。$\Omegac$ 就是\textbf{半功率点}。注意 $\Omegac$ 不一定等于通带截止频率 $\Omegap$——后面我们有得算。

    \item \textbf{「最大平坦」的数学含义：}
    在 $\Omega = 0$ 处，$|H(j\Omega)|^2$ 的前 $2N-1$ 阶导数全是 0。换句话说，它在直流附近「粘」在 $1$ 上，能粘多紧粘多紧。代价是\textbf{过渡带滚降不够快}——平坦的代价是慢。
\end{enumerate}

% --- 切比雪夫 ---
\subsubsection{切比雪夫幅度平方函数：用波纹换陡峭}

$$
\boxed{\absSq{H(j\Omega)} = \frac{1}{1 + \varepsilon^2\, C_N^2\!\left(\dfrac{\Omega}{\Omegap}\right)}}
$$

其中 $C_N(x) = \cos(N\arccos x)$（$|x|\le 1$）或 $\ch(N\arch x)$（$|x|>1$）。

\textbf{为什么长这样？算力经济学的博弈：}

\begin{enumerate}[leftmargin=2em]
    \item \textbf{$\varepsilon$（波纹因子）：} 决定通带内允许的波动幅度。$\varepsilon$ 越小，波纹越浅，滤波器越像巴特沃斯；$\varepsilon$ 越大，通带越「浪」，但过渡带也更陡。$0<\varepsilon<1$ 是典型取值。

    \item \textbf{切比雪夫多项式 $C_N(x)$ 的魔术：}
    在通带内（$|x|\le 1$），$C_N(x)$ 在 $[-1, 1]$ 之间等幅振荡 $N+1$ 次。这意味着 $|H|^2$ 在 $1$ 和 $1/(1+\varepsilon^2)$ 之间\textbf{等波纹地波动}。这不是缺点，而是刻意设计的——把通带误差均匀摊开，把宝贵的「容许误差」用到极致。

    \item \textbf{「等波纹逼近」的数学哲学：}
    切比雪夫证明了：在所有 $N$ 阶多项式中，$C_N(x)$ 在 $[-1,1]$ 上与零的最大偏差最小。翻译成工程语言：\textbf{给定通带容差，切比雪夫用相同的阶数在过渡带跑得最快}。用通带里的「可控波纹」为代价，换取更陡的过渡带——这就是滤波器设计中的「算力经济学」。

    \item \textbf{阻带的起飞：}
    一旦进入阻带（$|x|>1$），$C_N(x)$ 从 $\ch$ 函数接管，指数级增长。$|H|^2$ 被 $C_N^2$ 急速压向零，过渡带斜率远超同阶巴特沃斯。
\end{enumerate}

\begin{insight}
\textbf{一句话对比：}
\begin{itemize}
    \item 巴特沃斯 = 「我老老实实地平坦，别嫌我慢。」
    \item 切比雪夫 = 「我在通带里故意晃几下，但保证不晃出圈子，过渡带跑得飞快。」
\end{itemize}
选定哪个，取决于甲方更在乎通带平坦度还是过渡带陡峭度。
\end{insight}

% ---------------------------------------------------------------------------
\subsection{极点分布的几何图景：圆与椭圆的宿命}

现在进入最具美感的环节。幅度平方函数决定之后，我们通过
\[
|H(j\Omega)|^2 = H(s)H(-s)\big|_{s=j\Omega}
\]
把 $s$ 请回来。解 $1 + (\text{分母}) = 0$ 得到的根，画在 $s$ 平面上，就是极点。

% --- 巴特沃斯极点 ---
\subsubsection*{巴特沃斯：完美圆周上的等分点}

令 $s = j\Omega$，有 $|H(s)|^2 = H(s)H(-s) = \dfrac{1}{1 + \bigl(-s^2/\Omegac^2\bigr)^{N}}$。

极点是 $1 + \bigl(-s^2/\Omegac^2\bigr)^{N}=0$ 的解，即：

$$
\bigl(-s^2/\Omegac^2\bigr)^{N} = -1 = e^{j\pi(2k-1)}
$$

整理得（设 $N$ 为奇数或偶数，统一下标）：

\vspace{4pt}
\begin{center}
\fbox{%
\begin{minipage}{0.55\textwidth}
\centering\large
$\displaystyle s_k = \Omegac \cdot \exp\!\left[j\frac{\pi}{2}\left(1 + \frac{2k-1}{N}\right)\right]$\vspace{6pt}

$k = 1, 2, \dots, N$ （仅取左半平面极点）
\end{minipage}%
}
\end{center}
\vspace{4pt}

\begin{insight}
\textbf{关键发现：}
\begin{itemize}
    \item 所有极点均匀分布在以 $\Omegac$ 为半径的\textbf{圆周}上。
    \item 相邻极点之间的角度间隔为 $\pi/N$。
    \item 极点永远不会落在 $j\Omega$ 轴上（那会造成无穷大增益，滤波器不稳定）。
    \item 左半平面的 $N$ 个极点构成了因果稳定滤波器的 $H(s)$。
    \item 巴特沃斯极点 = 完全对称 = 无波纹 = 「秩序之美」。
\end{itemize}
\end{insight}

% --- 切比雪夫极点 ---
\subsubsection*{切比雪夫 I 型：椭圆上的分布}

切比雪夫滤波器的极点不在圆上，而是落在 $s$ 平面长轴在 $j\Omega$ 轴上、短轴在实轴上的\textbf{椭圆}上：

$$
s_k = -\Omegap \cdot \sh(\alpha)\,\sin\theta_k \;+\; j\,\Omegap \cdot \ch(\alpha)\,\cos\theta_k
$$

其中 $\alpha = \dfrac{1}{N}\,\arsh\!\left(\dfrac{1}{\varepsilon}\right)$，$\theta_k = \dfrac{\pi}{2}\cdot\dfrac{2k-1}{N}$。

\begin{insight}
\textbf{椭圆解读：}
\begin{itemize}
    \item 椭圆的短半轴 = $\Omegap \cdot \sh(\alpha)$（实轴方向，决定极点「阻尼」）
    \item 椭圆的长半轴 = $\Omegap \cdot \ch(\alpha)$（虚轴方向，决定极点「频率」）
    \item $\varepsilon \to 0$ 时，$\alpha \to \infty$，$\sh\alpha \approx \ch\alpha \approx e^\alpha/2$，椭圆退化为圆——切比雪夫退化为巴特沃斯。
    \item $\varepsilon$ 越大，椭圆越扁，极点越靠近 $j\Omega$ 轴，通带波纹越大，但过渡带越陡。
\end{itemize}
\end{insight}

% --- TikZ 图：圆 vs 椭圆并排 ---
\begin{figure}[htbp]
\centering
\begin{tikzpicture}[scale=0.85]

    % ===== 左：巴特沃斯 N=6 =====
    \begin{scope}[xshift=-4.0cm]
        % 坐标轴
        \draw[->, gray] (-3, 0) -- (3, 0) node[right] {$\sigma$};
        \draw[->, gray] (0, -3) -- (0, 3) node[above] {$j\Omega$};

        % 单位圆（归一化 Omega_c=1）
        \draw[butterworth, thick, dashed] (0,0) circle (2.2);

        % 左半平面极点 (k=1,2,3: angles 120°, 180°, 240°)
        \foreach \k in {1,2,3} {
            \pgfmathsetmacro{\ang}{90 + (2*\k-1)*180/6}
            \fill[butterworth] ({2.2*cos(\ang)}, {2.2*sin(\ang)}) circle (3pt);
            \draw[butterworth, very thin] (0,0) -- ({2.2*cos(\ang)}, {2.2*sin(\ang)});
        }
        % 右半平面极点（灰色, k=4,5,6: angles 300°, 0°, 60°）
        \foreach \k in {4,5,6} {
            \pgfmathsetmacro{\ang}{90 + (2*\k-1)*180/6}
            \fill[gray!50] ({2.2*cos(\ang)}, {2.2*sin(\ang)}) circle (2.5pt);
        }

        % 标注
        \node[butterworth, font=\footnotesize\bfseries] at (1.2, 0.8) {$30^\circ$};
        \node[butterworth, font=\small\bfseries] at (0, -2.8) {巴特沃斯 ($N=6$)};
        \node[font=\footnotesize] at (2.3, 1.4) {半径 $=\Omegac$};

        % σ 负方向指示（左半平面）
        \draw[<-, thick, butterworth!60] (-1.3, 0.5) -- (-2.0, 1.2);
        \node[butterworth!80, font=\scriptsize] at (-2.5, 1.5) {稳定极点};
    \end{scope}

    % ===== 右：切比雪夫 N=6 =====
    \begin{scope}[xshift=4.0cm]
        % 坐标轴
        \draw[->, gray] (-3, 0) -- (3, 0) node[right] {$\sigma$};
        \draw[->, gray] (0, -3) -- (0, 3) node[above] {$j\Omega$};

        % 椭圆参数
        % epsilon=0.5, N=6 -> alpha = arsinh(1/0.5)/6 = arsinh(2)/6 ≈ 1.4436/6 ≈ 0.2406
        % a = sin(real) cos(imag)? Let me use explicit values.
        % sh(alpha) = sh(0.2406) ≈ 0.243, ch(alpha) ≈ 1.029
        % semi-minor = Omega_p * sh ≈ 1 * 0.243 * scale
        % semi-major = Omega_p * ch ≈ 1 * 1.029 * scale
        % Let's scale for visibility
        \draw[chebyshev, thick, dashed] (0,0) ellipse (0.55 and 2.35);

        % 极点 N=6 (左半平面椭圆上)
        \foreach \k in {1,2,3,4,5,6} {
            \pgfmathsetmacro{\theta}{90*(2*\k-1)/6}
            \pgfmathsetmacro{\x}{-0.55*sin(\theta)}
            \pgfmathsetmacro{\y}{2.35*cos(\theta)}
            \fill[chebyshev] (\x, \y) circle (3pt);
        }
        % 右半平面极点（灰色）
        \foreach \k in {1,2,3,4,5,6} {
            \pgfmathsetmacro{\theta}{90*(2*\k-1)/6}
            \pgfmathsetmacro{\x}{0.55*sin(\theta)}
            \pgfmathsetmacro{\y}{2.35*cos(\theta)}
            \fill[gray!50] (\x, \y) circle (2.5pt);
        }

        \node[chebyshev, font=\small\bfseries] at (0, -2.8) {切比雪夫 I 型 ($N=6$)};
        \node[font=\footnotesize] at (1.5, 1.8) {长轴 $=\Omegap\ch\alpha$};
        \node[font=\footnotesize] at (1.2, 0.4) {短轴 $=\Omegap\sh\alpha$};

        % 稳定性指示
        \draw[<-, thick, chebyshev!60] (-0.8, 0.5) -- (-1.8, 1.3);
        \node[chebyshev!80, font=\scriptsize] at (-2.4, 1.6) {稳定极点};
    \end{scope}

\end{tikzpicture}
\caption{巴特沃斯（左）与切比雪夫~I 型（右）极点分布几何对比。
圆上均匀分布的极点带来「最大平坦」；椭圆分布则通过挤压换取通带波纹与过渡带陡度。}
\label{fig:pole_geometry}
\end{figure}

% --- 共轭极点合并 ---
\subsubsection*{共轭极点合并为实系数二阶节——化简为工程实现}

极点成对出现（共轭对），每一对共轭极点合并为一个实系数的二阶多项式：

\begin{align*}
\text{单极点（实轴）:}&\quad (s - s_k), \quad s_k \in \mathbb{R} \\
\text{共轭对:}&\quad (s - s_k)(s - s_k^*) = s^2 - 2\operatorname{Re}\{s_k\}s + |s_k|^2
\end{align*}

\begin{insight}
\textbf{工程意义：}
\textit{复数系数}的滤波器无法用运放实现。把共轭极点合并为实系数二阶节之后，每个二阶节恰好对应一个 Sallen-Key 或 MFB 二阶单元。级联后就是完整的滤波器硬件。这就是为什么所有归一化表格都用「一阶节 + 多个二阶节」的形式列出——表格里的每行就是一块运放电路。
\end{insight}

\newpage
% =============================================================================
%  PART II: 4步傻瓜设计法
% =============================================================================
\section{终极工程模板：模拟滤波器 4 步傻瓜设计法}

先把变量声明清楚：
\[
\begin{aligned}
\Omegap &= 2\pi f_p &&\text{—— 通带截止角频率} \\
\Omegas &= 2\pi f_s &&\text{—— 阻带截止角频率} \\
\alpha_p &\ \db &&\text{—— 通带最大允许衰减} \\
\alpha_s &\ \db &&\text{—— 阻带最小所需衰减}
\end{aligned}
\]

目标：找到阶数 $N$ 和 $3\db$ 截止频率 $\Omegac$（巴特沃斯）或波纹参数 $\varepsilon$（切比雪夫），进而写出 $H(s)$。

\bigskip
\hrule
\subsection*{Step 1：探底线——求最小阶数 $N$}

\textbf{巴特沃斯：}
\[
N \ge \frac{\lg\!\left(\dfrac{10^{\alpha_s/10}-1}{10^{\alpha_p/10}-1}\right)}
                {2\,\lg\!\left(\dfrac{\Omegas}{\Omegap}\right)}
\]

\textbf{切比雪夫 I 型：}
\[
N \ge \frac{\arch\!\left(\sqrt{\dfrac{10^{\alpha_s/10}-1}{10^{\alpha_p/10}-1}}\right)}
                {\arch\!\left(\dfrac{\Omegas}{\Omegap}\right)}
\]

\textbf{记忆口诀：} 巴特沃斯看「对数比」；切比雪夫看「反双曲余弦比」。两者的 $N$ 公式分母都是过渡带的「压缩比」$\Omegas/\Omegap$——过渡带越窄，$N$ 越大，很直观。

\textbf{切比雪夫 $\arch$ 的数值技巧：} $\arch(x) = \ln\!\bigl(x + \sqrt{x^2-1}\,\bigr)$，$x>1$。考场没有 $\arch$ 按钮就用这个。

\bigskip
\hrule
\subsection*{Step 2：定硬件——阶数向上取整}

\[
N_{\text{实际}} = \ceil{N_{\text{计算}}}
\]

\textbf{工程师的直觉：} 滤波器的阶数必须是整数——你不可能做 $3.2$ 个运放。取整意味着「硬件冗余」：
\begin{itemize}
    \item 多出来的 $0.x$ 阶不会浪费，它会转化为额外的衰减余量。
    \item 这个余量我们可以「分配」给通带（把截止频率往外挪）或者吃掉（保持原来指标更轻松地满足）。
\end{itemize}

\bigskip
\hrule
\subsection*{Step 3：分蛋糕——计算真实参数}

\textbf{巴特沃斯——两种策略：}

\begin{minipage}[t]{0.47\textwidth}
\centering
\textbf{策略 A：「保通带」}\\[4pt]
\fbox{$\displaystyle \Omegac = \frac{\Omegap}{\bigl(10^{\alpha_p/10}-1\bigr)^{1/(2N)}}$}\\[8pt]
$3\db$ 点取在通带指标刚好满足的位置。阻带自动获得额外余量。
\end{minipage}
\hfill
\begin{minipage}[t]{0.47\textwidth}
\centering
\textbf{策略 B：「保阻带」}\\[4pt]
\fbox{$\displaystyle \Omegac = \frac{\Omegas}{\bigl(10^{\alpha_s/10}-1\bigr)^{1/(2N)}}$}\\[8pt]
$3\db$ 点取在阻带指标刚好满足的位置。通带自动获得额外余量。
\end{minipage}

\vspace{8pt}

\textbf{切比雪夫 I 型——波纹因子：}
\[
\varepsilon = \sqrt{10^{\alpha_p/10} - 1}
\]

\textbf{选哪个策略？} 如果甲方对通带平坦度极端敏感（比如音频），选「保通带」；如果甲方更在意阻带衰减（比如抗混叠），选「保阻带」。如果没特别说明，一般保通带。

\bigskip
\hrule
\subsection*{Step 4：查表与施工——去归一化}

\begin{enumerate}[leftmargin=1.5em, label=\textbf{\arabic*.}]
    \item \textbf{查归一化表：} 根据 $N$ 查得归一化传递函数 $\HanNorm$（$\Omegac=1$ 或 $\Omegap=1$ 的原型）。
    \item \textbf{去归一化：}
    \begin{itemize}
        \item \textbf{巴特沃斯：} 令 $s \to s/\Omegac$，即
        \[
        \Han = \HanNorm\big|_{s \to s/\Omegac}
        \]
        \item \textbf{切比雪夫：} 令 $s \to s/\Omegap$，即
        \[
        \Han = \HanNorm\big|_{s \to s/\Omegap}
        \]
    \end{itemize}
    \item \textbf{化简：} 将替换后的表达式展开为实系数多项式之比，即
    \[
    \Han = \frac{b_0 + b_1 s + \cdots}{a_0 + a_1 s + \cdots + a_N s^N}
    \]
\end{enumerate}

\begin{pitfall}
去归一化的核心陷阱：巴特沃斯用 $\Omegac$ 归一化，切比雪夫用 $\Omegap$ 归一化。两者查的表不同，换的变量也不同——千万别在考场上张冠李戴！
\end{pitfall}

\newpage
% =============================================================================
%  APPENDIX: 归一化表
% =============================================================================
\subsection*{附：归一化低通原型多项式表（简表）}

以下给出巴特沃斯和切比雪夫（$1\db$ 波纹）的归一化传递函数分母多项式 $D_{an}(s)$，分子均为 $1$。

\vspace{4pt}

\noindent
\begin{minipage}{\textwidth}
\centering\small
\textbf{巴特沃斯}（$\Omegac=1$）\\[4pt]
\setlength{\tabcolsep}{6pt}
\begin{tabular}{c|p{0.75\textwidth}}
\toprule
$N$ & $D_{an}(s)$ 因式分解形式 \\
\midrule
1  & $s+1$ \\
2  & $s^2 + 1.4142s + 1$ \\
3  & $(s+1)(s^2 + s + 1)$ \\
4  & $(s^2 + 0.7654s + 1)(s^2 + 1.8478s + 1)$ \\
5  & $(s+1)(s^2 + 0.6180s + 1)(s^2 + 1.6180s + 1)$ \\
6  & $(s^2 + 0.5176s + 1)(s^2 + 1.4142s + 1)(s^2 + 1.9319s + 1)$ \\
\bottomrule
\end{tabular}
\end{minipage}

\vspace{10pt}

\noindent
\begin{minipage}{\textwidth}
\centering\small
\textbf{切比雪夫 I 型}（$1\db$ 波纹，$\Omegap=1$）\\[4pt]
\setlength{\tabcolsep}{6pt}
\begin{tabular}{c|p{0.75\textwidth}}
\toprule
$N$ & $D_{an}(s)$ 因式分解形式 \\
\midrule
1  & $s+1.9652$ \\
2  & $s^2 + 1.0977s + 1.1025$ \\
3  & $(s+0.4942)(s^2 + 0.4942s + 0.9942)$ \\
4  & $(s^2+0.2791s+0.9865)(s^2+0.6737s+0.2794)$ \\
5  & $(s+0.2895)(s^2+0.1790s+0.9883)(s^2+0.4684s+0.4293)$ \\
\bottomrule
\end{tabular}
\end{minipage}

\vspace{8pt}
\begin{pitfall}
上表仅为最常用的 $1\db$ 波纹切比雪夫。考场上请根据题目给出的表格为准。巴特沃斯表注意区分归一化到 $\Omegac$ 还是 $\Omegap$——大多数教材给出归一化到 $3\db$ 点的表格，与本文一致。
\end{pitfall}

\newpage
% =============================================================================
%  PART III: 经典精讲例题
% =============================================================================
\section{经典精讲例题}

% ==============================
%  例题 1: 巴特沃斯
% ==============================
\subsection{例题 1：巴特沃斯低通滤波器设计（保通带）}

\begin{tcolorbox}[enhanced, colback=white, colframe=butterworth!60,
    title={\textbf{题目}}, fonttitle=\bfseries, sharp corners, breakable]

\textbf{设计指标：}
\[
f_p = 1\ \mathrm{kHz},\quad \alpha_p = 1\db,\qquad
f_s = 3\ \mathrm{kHz},\quad \alpha_s = 40\db
\]

用巴特沃斯逼近设计模拟低通滤波器，求：
\begin{enumerate}[label=(\arabic*)]
    \item 所需最小阶数 $N$；
    \item 采用「保通带」策略的 $3\db$ 截止频率 $f_c$；
    \item 归一化传递函数 $\HanNorm$；
    \item 去归一化后的实际传递函数 $\Han$（化简为实系数多项式之比）。
\end{enumerate}
\end{tcolorbox}

\vspace{8pt}

% ----- 详细解答 -----
\begin{center}
    \textbf{——— 详细解析 ———}
\end{center}

\textbf{\color{butterworth!90}Step 1: 探底线（求 $N$）}

\begin{align*}
\Omega_p &= 2\pi \times 1000 = 2000\pi \ \rads \\
\Omega_s &= 2\pi \times 3000 = 6000\pi \ \rads \\[4pt]
N &\ge \frac{\lg\!\left(\dfrac{10^{40/10}-1}{10^{1/10}-1}\right)}
                {2\,\lg\!\left(\dfrac{6000\pi}{2000\pi}\right)}
     = \frac{\lg\!\left(\dfrac{10^4-1}{10^{0.1}-1}\right)}
                {2\,\lg(3)} \\[4pt]
  &= \frac{\lg\!\left(\dfrac{9999}{0.2589}\right)}{2 \times 0.4771}
     = \frac{\lg(38620)}{0.9542}
     = \frac{4.5868}{0.9542}
     = 4.807
\end{align*}

\begin{insight}
注意 $10^{0.1} \approx 1.2589$，$10^{0.1}-1 \approx 0.2589$。$1\db$ 这个数在滤波器设计里极其常见——它对应的功率比约为 $1.259$，幅度比约为 $1.122$。记住 $10^{0.1}\approx 1.259$ 可以在考场上省去按计算器的时间。
\end{insight}

\textbf{\color{deepgreen!90}Step 2: 定硬件（取整）}

\[
N = \ceil{4.807} = \mathbf{5}
\]

取整到 $5$ 阶，意味着「硬件冗余」约为 $0.193$ 阶——这部分余量将被 Step 3 消化。

\vspace{4pt}

\textbf{\color{chebyshev!90}Step 3: 分蛋糕（计算 $\Omega_c$）}

甲方要求「保通带」—— 通带指标必须精确满足，阻带多吃点余量无妨：

\begin{align*}
\Omegac &= \frac{\Omega_p}{\bigl(10^{\alpha_p/10}-1\bigr)^{1/(2N)}}
         = \frac{2000\pi}{(10^{0.1}-1)^{1/10}} \\[4pt]
        &= \frac{2000\pi}{(0.2589)^{0.1}}
         = \frac{2000\pi}{0.8734} \\[4pt]
        &\approx 2289.8\pi \ \rads \\[4pt]
f_c &= \frac{\Omega_c}{2\pi} \approx 1144.9 \ \mathrm{Hz}
\end{align*}

\begin{insight}
注意 $\Omega_c > \Omega_p$（$1145\Hz > 1000\Hz$）——取整后多出的硬件余量被我们「分」到 $3\db$ 点上了：把截止频率往外推，通带指标依然安然无恙。如果取整后仍然卡在原来的 $\Omega_c$ 上，那阻带余量会大到离谱——那是浪费。
\end{insight}

\textbf{\color{black!80}Step 4: 查表与施工（去归一化）}

查巴特沃斯归一化表，$N=5$：

\[
\HanNorm = \frac{1}{(s+1)(s^2 + 0.6180s + 1)(s^2 + 1.6180s + 1)}
\]

去归一化：令 $s \to s/\Omegac = s/(2289.8\pi)$。

\textbf{—— 逐项替换 ——}

\begin{align*}
\text{一阶因子:}&\quad (s/\Omegac + 1)
    = \frac{1}{\Omegac}s + 1
    = \frac{s + \Omegac}{\Omegac} \\[4pt]
\text{二阶因子 (a):}&\quad (s/\Omegac)^2 + 0.6180(s/\Omegac) + 1
    = \frac{s^2 + 0.6180\,\Omegac\,s + \Omegac^2}{\Omegac^2} \\[4pt]
\text{二阶因子 (b):}&\quad (s/\Omegac)^2 + 1.6180(s/\Omegac) + 1
    = \frac{s^2 + 1.6180\,\Omegac\,s + \Omegac^2}{\Omegac^2}
\end{align*}

\textbf{—— 合并 ——}

\begin{align*}
\Han &= \frac{1}
    {\dfrac{s + \Omegac}{\Omegac}
     \cdot \dfrac{s^2 + 0.6180\,\Omegac\,s + \Omegac^2}{\Omegac^2}
     \cdot \dfrac{s^2 + 1.6180\,\Omegac\,s + \Omegac^2}{\Omegac^2}} \\[8pt]
    &= \frac{\Omegac^5}
       {(s + \Omegac)\,
        (s^2 + 0.6180\,\Omegac\,s + \Omegac^2)\,
        (s^2 + 1.6180\,\Omegac\,s + \Omegac^2)}
\end{align*}

\begin{pitfall}
去归一化之后分子的 $\Omegac^N$ 因子\textbf{不能丢}！分子 $\Omegac^5$ 保证了直流增益 $|H(j0)| = 1$。如果忘了写这个因子，$H(0) \neq 1$，滤波器会把直流信号也衰减掉——这在绝大多数应用中都是不可接受的。
\end{pitfall}

\textbf{—— 代入数值 }$\Omegac \approx 2289.8\pi = 7193.3\ \rads$：

\[
\boxed{\displaystyle
\Han = \frac{(7193.3)^5}
    {(s + 7193.3)\,
    (s^2 + 4445.5\,s + 7193.3^2)\,
    (s^2 + 11638.8\,s + 7193.3^2)}
}
\]

%（可进一步乘开得到标准多项式形式，但考试一般要求到因式分解形式即可。）


\newpage
% ==============================
%  例题 2: 切比雪夫
% ==============================
\subsection{例题 2：切比雪夫 I 型低通滤波器设计}

\begin{tcolorbox}[enhanced, colback=white, colframe=chebyshev!60,
    title={\textbf{题目}}, fonttitle=\bfseries, sharp corners, breakable]

\textbf{设计指标：}
\[
f_p = 2\ \mathrm{kHz},\quad \alpha_p = 1\db,\qquad
f_s = 3.5\ \mathrm{kHz},\quad \alpha_s = 30\db
\]

用切比雪夫 I 型逼近设计模拟低通滤波器，求：
\begin{enumerate}[label=(\arabic*)]
    \item 波纹因子 $\varepsilon$；
    \item 所需最小阶数 $N$；
    \item 归一化传递函数 $\HanNorm$；
    \item 去归一化后的实际传递函数 $\Han$。
\end{enumerate}
\end{tcolorbox}

\vspace{8pt}

% ----- 详细解答 -----
\begin{center}
    \textbf{——— 详细解析 ———}
\end{center}

\textbf{\color{butterworth!90}Step 1a: 计算波纹因子}

\[
\varepsilon = \sqrt{10^{\alpha_p/10} - 1} = \sqrt{10^{0.1} - 1}
            = \sqrt{1.2589 - 1} = \sqrt{0.2589} \approx 0.5088
\]

物理含义：$\varepsilon = 0.5088$ 意味着通带内 $|H|$ 在 $1$ 和 $1/\sqrt{1+\varepsilon^2} \approx 1/\sqrt{1.2589} \approx 0.8913$（恰好 $-1\db$）之间等幅振荡——波纹幅度恰好等于甲方规定的 $\alpha_p$。

\vspace{4pt}

\textbf{\color{butterworth!90}Step 1b: 探底线（求 $N$）}

\begin{align*}
\Omega_p &= 4000\pi\ \rads,\qquad \Omega_s = 7000\pi\ \rads \\[4pt]
N &\ge \frac{\arch\!\left(\sqrt{\dfrac{10^{30/10}-1}{10^{1/10}-1}}\right)}
                {\arch(3.5/2)}
     = \frac{\arch\!\left(\sqrt{\dfrac{1000-1}{0.2589}}\right)}
                {\arch(1.75)} \\[4pt]
  &= \frac{\arch(62.12)}{\arch(1.75)}
\end{align*}

使用 $\arch(x) = \ln\bigl(x + \sqrt{x^2-1}\bigr)$：
\begin{align*}
\arch(62.12) &= \ln\!\bigl(62.12 + \sqrt{62.12^2-1}\bigr)
              = \ln(62.12 + 62.11) = \ln(124.23) = 4.822 \\[2pt]
\arch(1.75)  &= \ln\!\bigl(1.75 + \sqrt{1.75^2-1}\bigr)
              = \ln(1.75 + 1.4361) = \ln(3.1861) = 1.159 \\[4pt]
N &\ge \frac{4.822}{1.159} = 4.160
\end{align*}

\begin{insight}
\textbf{对比意识：} 如果这是巴特沃斯，$N$ 会是多少？带入巴特沃斯公式：
$N_{\text{BW}} \ge \lg(3859)/(2\lg1.75) = 3.587/(2\times0.243)=7.38$，取整后 $N=8$。
切比雪夫只用 $N=5$ 阶就达到了同样的阻带衰减——\textbf{省了 $3$ 个运放！}这就是「用波纹换阶数」的工程价值。
\end{insight}

\textbf{\color{deepgreen!90}Step 2: 定硬件（取整）}

\[
N = \ceil{4.160} = \mathbf{5}
\]

\vspace{4pt}

\textbf{\color{chebyshev!90}Step 3: 参数确认}

切比雪夫不需要像巴特沃斯那样重算截止频率——$\Omegap$ 本身就是归一化参考点，取整后通带指标自然满足（因为 $N$ 增大只会让通带波纹更多但不超出 $\varepsilon$ 的限制）。

实际参数：
\[
\varepsilon = 0.5088,\quad N = 5,\quad \Omega_p = 4000\pi\ \rads
\]

\textbf{\color{black!80}Step 4: 查表与施工（去归一化）}

查切比雪夫 I 型归一化表（$1\db$ 波纹），$N=5$：

\[
\HanNorm = \frac{1}{(s+0.2895)(s^2+0.1790s+0.9883)(s^2+0.4684s+0.4293)}
\]

去归一化：令 $s \to s/\Omegap = s/(4000\pi)$。

\textbf{—— 逐项替换 ——}

\begin{align*}
\text{一阶因子:}&\quad s/\Omegap + 0.2895
    = \frac{s + 0.2895\Omegap}{\Omegap} \\[4pt]
\text{二阶因子 (a):}&\quad (s/\Omegap)^2 + 0.1790(s/\Omegap) + 0.9883
    = \frac{s^2 + 0.1790\Omegap\,s + 0.9883\Omegap^2}{\Omegap^2} \\[4pt]
\text{二阶因子 (b):}&\quad (s/\Omegap)^2 + 0.4684(s/\Omegap) + 0.4293
    = \frac{s^2 + 0.4684\Omegap\,s + 0.4293\Omegap^2}{\Omegap^2}
\end{align*}

\textbf{—— 合并 ——}

\begin{align*}
\Han &= \frac{1}
    {\dfrac{s + 0.2895\Omegap}{\Omegap}
     \cdot \dfrac{s^2 + 0.1790\Omegap\,s + 0.9883\Omegap^2}{\Omegap^2}
     \cdot \dfrac{s^2 + 0.4684\Omegap\,s + 0.4293\Omegap^2}{\Omegap^2}} \\[8pt]
    &= \frac{\Omegap^5}
       {(s + 0.2895\Omegap)\,
        (s^2 + 0.1790\Omegap\,s + 0.9883\Omegap^2)\,
        (s^2 + 0.4684\Omegap\,s + 0.4293\Omegap^2)}
\end{align*}

代入 $\Omegap = 4000\pi \approx 12566.4\ \rads$：

\[
\boxed{\displaystyle
\Han = \frac{(12566.4)^5}
    {(s + 3638.0)\,
    (s^2 + 2249.4\,s + 0.9883\cdot\Omegap^2)\,
    (s^2 + 5886.4\,s + 0.4293\cdot\Omegap^2)}
}
\]

\begin{pitfall}
切比雪夫去归一化时，\textbf{归一化表中的常数系数在替换后保持不变}（如 $0.2895$ 直接乘以 $\Omegap$，$0.9883$ 乘以 $\Omegap^2$），没有额外的功率补偿因子。这跟巴特沃斯不同——巴特沃斯去归一化后所有系数恰好凑成 $\Omegac$ 的整数次幂，而切比雪夫的系数保留了归一化表中的小数。
\end{pitfall}


\newpage
% =============================================================================
%  PART IV: 肌肉记忆训练营
% =============================================================================
\section{肌肉记忆训练营（刷题区）}

\textbf{使用说明：}
以下 4 道题目完全覆盖巴特沃斯和切比雪夫滤波器的常见考法。先独立完成，再对照答案核对核心节点数据。不要看详细解析——你对的是\textbf{数字}，不是思路。思路要在前面两节里形成。

\vspace{10pt}

% ---- 练习 1: 巴特沃斯（保通带） ----
\begin{tcolorbox}[enhanced, colback=white, colframe=black!40,
    title={\textbf{练习 1：巴特沃斯·保通带}}, fonttitle=\bfseries, sharp corners, breakable]

\textbf{指标：} $f_p = 800\Hz$，$\alpha_p = 2\db$；$f_s = 2400\Hz$，$\alpha_s = 30\db$。

求 $N$，用「保通带」策略求 $f_c$，写出最终的 $H(s)$ 因式分解形式（归一化查表，不展开多项式）。
\end{tcolorbox}

\vspace{4pt}
\begin{answerbox}
$N = 4$；$f_c \approx 854.8\Hz$；$\Omegac \approx 5371.2\rads$。

\[
\HanNorm = \frac{1}{(s^2+0.7654s+1)(s^2+1.8478s+1)}
\]

\[
\Han = \frac{\Omegac^4}
    {(s^2 + 0.7654\Omegac\,s + \Omegac^2)\,
     (s^2 + 1.8478\Omegac\,s + \Omegac^2)}
\]
\end{answerbox}

\vspace{16pt}

% ---- 练习 2: 巴特沃斯（保阻带） ----
\begin{tcolorbox}[enhanced, colback=white, colframe=black!40,
    title={\textbf{练习 2：巴特沃斯·保阻带}}, fonttitle=\bfseries, sharp corners, breakable]

\textbf{指标：} $f_p = 500\Hz$，$\alpha_p = 1\db$；$f_s = 1200\Hz$，$\alpha_s = 50\db$。

求 $N$，用「保阻带」策略求 $f_c$，写出 $\Han$ 的因式分解形式。
\end{tcolorbox}

\vspace{4pt}
\begin{answerbox}
$N_{\text{计算}} \approx 8.60$，取 $N = 9$。

$\Omegac$（保阻带）$\approx 3052.6\rads$，$f_c \approx 485.9\Hz$。

\[
\HanNorm = \frac{1}{(s+1)(s^2+0.3473s+1)(s^2+s+1)(s^2+1.5321s+1)(s^2+1.8794s+1)}
\]

\[
\Han = \HanNorm\big|_{s\to s/\Omegac}\quad\text{（代入 $\Omegac$，分子配 $\Omegac^9$）}
\]
\end{answerbox}

\vspace{16pt}

% ---- 练习 3: 切比雪夫 ----
\begin{tcolorbox}[enhanced, colback=white, colframe=black!40,
    title={\textbf{练习 3：切比雪夫 I 型·标准设计}}, fonttitle=\bfseries, sharp corners, breakable]

\textbf{指标：} $f_p = 1500\Hz$，$\alpha_p = 0.5\db$；$f_s = 2800\Hz$，$\alpha_s = 45\db$。

求 $\varepsilon$、$N$，写出 $\Han$ 因式分解形式（归一化参数可用字母表示）。
\end{tcolorbox}

\vspace{4pt}
\begin{answerbox}
$\varepsilon = \sqrt{10^{0.05}-1} \approx 0.3493$。

$\arch$ 参数：$\arch(1.8667) \approx 1.2146$；$\arch(\sqrt{31621.7/0.1220}) \approx \arch(509.1) \approx 6.926$。

$N_{\text{计算}} \approx 5.70$，取 $N = 6$。

查 $0.5\db$ 波纹表 $N=6$（本文仅附 $1\db$ 表，考场依题给表为准），得归一化 $\HanNorm$。

$\Omegap = 3000\pi$，去归一化：$s \to s/\Omegap$，分子配 $\Omegap^6$。
\[
\Han = \frac{\Omegap^6}{\prod_{k=1}^{3}(s^2 + a_k\Omegap s + b_k\Omegap^2)}
\]
（$a_k, b_k$ 查考场提供的 $0.5\db$ 归一化表）
\end{answerbox}

\vspace{16pt}

% ---- 练习 4: 切比雪夫·大衰减 ----
\begin{tcolorbox}[enhanced, colback=white, colframe=black!40,
    title={\textbf{练习 4：切比雪夫 I 型·大衰减考验}}, fonttitle=\bfseries, sharp corners, breakable]

\textbf{指标：} $f_p = 100\Hz$，$\alpha_p = 2\db$；$f_s = 210\Hz$，$\alpha_s = 60\db$。

求 $\varepsilon$、$N$，简要写出去归一化后的 $\Han$。
\end{tcolorbox}

\vspace{4pt}
\begin{answerbox}
$\varepsilon = \sqrt{10^{0.2}-1} \approx 0.7648$。

$\arch$ 分子：$\sqrt{(10^6-1)/(10^{0.2}-1)} \approx \sqrt{999999/0.5849} \approx \sqrt{1.709\times10^6} \approx 1307.3$。

$\arch(1307.3) \approx 7.869$。

$\arch(\Omegas/\Omegap) = \arch(210/100) = \arch(2.1) \approx 1.426$。

$N_{\text{计算}} \approx 5.52$，取 $N = 6$。

$\Omegap = 200\pi \approx 628.32\rads$。

查表 $N=6$（$1\db$ 波纹同一表），去归一化 $s \to s/\Omegap$，乘 $\Omegap^6$ 得 $\Han$。
\end{answerbox}

\vspace{20pt}

\begin{center}
\rule{0.6\textwidth}{0.5pt}\\[8pt]
{\large\kaishu\bfseries —— 祝刷题愉快，考场神挡杀神 ——}
\end{center}

% =============================================================================
%  APPENDIX
% =============================================================================
\newpage
\section*{附录 A：核心公式速查卡}

\begin{tcolorbox}[enhanced, colback=white, colframe=black!60,
    title={\textbf{公式速查：一张 A4 纸搞定滤波器设计}}, fonttitle=\bfseries\large,
    sharp corners, breakable]

\subsubsection*{巴特沃斯}
\[
|H(j\Omega)|^2 = \frac{1}{1+(\Omega/\Omegac)^{2N}},\qquad
\Omegac = \Omega_p\bigl(10^{\alpha_p/10}-1\bigr)^{-1/(2N)}
\]
\[
N = \frac{\lg\!\left(\dfrac{10^{\alpha_s/10}-1}{10^{\alpha_p/10}-1}\right)}{2\lg(\Omegas/\Omegap)},\qquad
s_k = \Omegac\cdot e^{j\frac{\pi}{2}\left(1+\frac{2k-1}{N}\right)}
\]

\subsubsection*{切比雪夫 I 型}
\[
|H(j\Omega)|^2 = \frac{1}{1+\varepsilon^2 C_N^2(\Omega/\Omegap)},\qquad
\varepsilon = \sqrt{10^{\alpha_p/10}-1}
\]
\[
N = \frac{\arch\!\left(\sqrt{\dfrac{10^{\alpha_s/10}-1}{10^{\alpha_p/10}-1}}\right)}{\arch(\Omegas/\Omegap)},\qquad
\arch(x) = \ln(x+\sqrt{x^2-1})
\]

\subsubsection*{去归一化}
\[
\text{巴特沃斯：}\ s\to s/\Omegac;\qquad
\text{切比雪夫：}\ s\to s/\Omegap
\]

\subsubsection*{角频率与频率}
\[
\Omega = 2\pi f,\qquad
\Omegap = 2\pi f_p,\qquad
\Omegas = 2\pi f_s,\qquad
\Omegac = 2\pi f_c
\]

\end{tcolorbox}

% =============================================================================
\end{document}
